%% kapa: exam=1 round=35 year=2024 period=1 form=A subject=real_estate subject_ko=부동산학원론 seq=6 number=86
%% source: Q-Net 감정평가사 기출문제 (cst00309) artlSeq=5213229 file=raw/1008/5213229/1교시 문제지 원본.pdf
%% conversion: kordoc 4.24.51 -> tools/gen1.py (auto); review=proofread-sonnet; flags=none
\begin{question}{86}
토지의 분류 및 용어에 관한 설명으로 옳은 것을 모두 고른 것은?
\begin{framedtext}
ㄱ. 획지(劃地)는 인위적, 자연적, 행정적 조건에 따라 다른 토지와 구별되는 가격수준이 비슷한 일단의 토지를 말한다. \par
ㄴ. 후보지(候補地)는 용도적 지역의 분류 중 세분된 지역 내에서 용도에 따라 전환되는 토지를 말한다. \par
ㄷ. 공지(空地)는 관련법령이 정하는 바에 따라 안전이나 양호한 생활환경을 확보하기 위해 건축하면서 남겨놓은 일정 면적의 토지를 말한다. \par
ㄹ. 갱지(更地)는 택지 등 다른 용도로 조성되기 이전 상태의 토지를 말한다. \par
\end{framedtext}
\begin{choices}
  \item ㄱ
  \item ㄹ
  \item ㄱ, ㄷ
  \item ㄴ, ㄹ
  \item ㄱ, ㄷ, ㄹ
\end{choices}
\end{question}
\answer{3}
